% =================================================
% Учительская версия рабочего листа
% =================================================

\worksheettitle
  {Точки через равные промежутки}
  {Учительская версия: ответы и методические комментарии}

\begin{teacherbox}[Цель листа]
  Сформировать различение количества точек и количества промежутков;
  научить находить расстояние между равноудалёнными точками, решать обратные
  задачи и переносить линейную модель на столбы, деревья и другие объекты.
\end{teacherbox}

\begin{checkpointbox}[Ответ к входной разминке]
  Пять точек образуют четыре промежутка. Желательно, чтобы ученик явно
  показал на рисунке четыре перехода между соседними точками.
\end{checkpointbox}

\section{Точки и промежутки}

\pointsvsintervalsfigure

\begin{propertybox}[Заполненный вывод]
  Между пятью точками находится \textbf{четыре} промежутка. Вообще,
  промежутков на \textbf{один} меньше, чем подряд отмеченных точек.
\end{propertybox}

\begin{practicebox}[Ответы к заданию 1]
  \begin{enumerate}[label=\alph*),itemsep=0.45em]
    \item \(9-3=6\) промежутков.
    \item \(27-8=19\) промежутков.
    \item \(16-15=1\) промежуток.
    \item \(56-41=15\) промежутков.
  \end{enumerate}
\end{practicebox}

\studentcountstrip

\begin{practicebox}[Ответы к заданию 2]
  \begin{enumerate}[label=\alph*),itemsep=0.45em]
    \item Шесть точек.
    \item \(7-2=5\) промежутков.
    \item \(5\cdot3=15\) см.
  \end{enumerate}
\end{practicebox}

\begin{teacherbox}[Методический акцент]
  Не стоит начинать с формулы. Сначала ученик должен увидеть переходы на
  коротком ряду: от 2-й до 3-й --- один, до 4-й --- два и т. д. После этого
  разность номеров воспринимается как число переходов, а не как правило для
  механического запоминания.
\end{teacherbox}

\section{Находим расстояние}

\begin{propertybox}[Заполненный алгоритм]
  Сначала найдите число \textbf{промежутков}, затем умножьте его на
  \textbf{длину одного промежутка}.
\end{propertybox}

\begin{practicebox}[Ответы к заданию 3]
  \begin{enumerate}[label=\alph*),itemsep=0.5em]
    \item \((12-5)\cdot7=49\) мм.
    \item \((27-8)\cdot7=133\) мм.
    \item \((44-31)\cdot7=91\) мм.
  \end{enumerate}
\end{practicebox}

\clearpage
\begin{thinkbox}[Ответ к заданию 4]
  Между 6-й и 14-й точками не восемь точек и не восемь промежутков, а
  \[
    14-6=8
  \]
  промежутков. В данном примере числовой ответ ученика случайно получился
  верным: \(8\cdot5=40\) см, но объяснение неверно. Это полезный пример:
  правильный ответ ещё не гарантирует правильного рассуждения.
\end{thinkbox}

\begin{teacherbox}[Тонкость формулировки]
  Между 6-й и 14-й точками находятся семь внутренних точек
  (7-я, 8-я, ..., 13-я), а если считать обе крайние, то всего девять точек.
  Промежутков при этом восемь. Просите ученика уточнять, считает ли он
  внутренние точки, все точки или промежутки.
\end{teacherbox}

\section{Обратные задачи}

\reverseexamplefigure

\begin{practicebox}[Ответы к заданию 5]
  \begin{enumerate}[label=\alph*),itemsep=0.55em]
    \item \(7-3=4\) промежутка; \(24:4=6\) см.
    \item \(19-11=8\) промежутков; \(56:8=7\) мм.
  \end{enumerate}
\end{practicebox}

\unknownpointfigure

\begin{practicebox}[Ответы к заданию 6]
  \begin{enumerate}[label=\alph*),itemsep=0.5em]
    \item Последняя точка имеет номер 17.
    \item \(28:4=7\) промежутков; \(12+7=19\).
    \item \(36:4=9\) промежутков; \(30-9=21\).
  \end{enumerate}
\end{practicebox}

\begin{teacherbox}[Диагностика направления]
  Если ученик всегда прибавляет число промежутков, попросите его показать
  направление движения на схеме. При движении вправо номера в принятом
  рисунке увеличиваются, при движении влево --- уменьшаются. Это свойство
  рисунка нужно проговорить; оно не следует только из расстояния.
\end{teacherbox}

\section{Столбы и деревья}

\postrowfigure

\begin{practicebox}[Ответ к заданию 7]
  \begin{enumerate}[label=\alph*),itemsep=0.45em]
    \item \(7-1=6\) промежутков.
    \item \(6\cdot4=24\) м.
  \end{enumerate}
\end{practicebox}

\begin{thinkbox}[Ответ к заданию 8]
  \[
    54:6=9
  \]
  промежутков. Так как фонари стоят в обоих концах ряда, фонарей на один
  больше:
  \[
    9+1=10.
  \]
\end{thinkbox}

\begin{teacherbox}[Граница применимости правила \(n-1\)]
  Правило работает, когда первый и последний объекты действительно входят в
  ряд. В задачах может быть сказано «между двумя столбами поставили ещё...»
  или «на участке длиной ... через каждые ...». Тогда нужно отдельно выяснить,
  есть ли объект в каждом конце. Без этого автоматическое прибавление единицы
  опасно.
\end{teacherbox}

\section{Чередование}

\alternatingtreesfigure

\begin{practicebox}[Ответы к заданию 9]
  \begin{enumerate}[label=\alph*),itemsep=0.45em]
    \item 17-е дерево --- сосна, так как 17 --- нечётное число.
    \item 28-е дерево --- берёза, так как 28 --- чётное число.
    \item Среди первых 20 деревьев 10 берёз.
    \item Среди первых 21 дерева 11 сосен.
  \end{enumerate}
\end{practicebox}

\begin{thinkbox}[Ответ к заданию 10]
  Между 7-м и 18-м деревьями
  \[
    18-7=11
  \]
  промежутков, поэтому расстояние равно
  \[
    11\cdot3=33\text{ м}.
  \]
  7-е дерево --- сосна, 18-е --- берёза.
\end{thinkbox}

\begin{checkpointbox}[Ответы к заданию 11]
  \begin{enumerate}[label=\alph*),itemsep=0.4em]
    \item \(-\): промежутков \(10-1=9\).
    \item \(+\): \(32-25=7\).
    \item \(+\).
    \item \(-\): сосны стоят на нечётных местах.
  \end{enumerate}
\end{checkpointbox}

\section{Типичные ошибки}

\begin{center}
  \renewcommand{\arraystretch}{1.42}
  \begin{tabularx}{\textwidth}{
    >{\raggedright\arraybackslash}p{0.29\textwidth}
    >{\raggedright\arraybackslash}p{0.32\textwidth}
    >{\raggedright\arraybackslash}X}
    \toprule
    \rowcolor{TableHeader}
    Ошибка & Что не понято & Наводящий вопрос \\
    \midrule
    Число промежутков равно номеру последней точки
      & Не учтён номер начальной точки
      & Сколько переходов от 8-й до 9-й? А до 10-й? \\
    \addlinespace[0.2em]
    Для \(n\) объектов используется \(n\) промежутков
      & Смешаны объекты и расстояния между ними
      & Сколько промежутков между двумя столбами? \\
    \addlinespace[0.2em]
    В обратной задаче умножают вместо деления
      & Не установлено, что общее состоит из равных частей
      & На сколько равных промежутков разделено расстояние? \\
    \addlinespace[0.2em]
    Вид дерева определяется без учёта первого объекта
      & Не замечена начальная позиция чередования
      & Что стоит на 1-м, 2-м, 3-м местах? \\
    \bottomrule
  \end{tabularx}
\end{center}

\begin{teacherbox}[Итоговая формулировка ученика]
  «Число промежутков между точками равно разности их номеров. Чтобы найти
  расстояние, число промежутков умножают на длину одного промежутка».
\end{teacherbox}
